% !TEX TS-program = xelatex 
% !TEX spellcheck = uk_UA
% !TEX options = --output-directory=build  
% !TEX encoding = UTF-8  

% *********** Приклад оформлення лабораторних робіт ************
% ************* Зміни Косенка О. та Шкарупило М.  **************
% ********************** На основі файлу ***********************
% ********** Приклад оформлення пояснювальної записки **********
% *********  до атестаційної роботи ступеня бакалавра **********
% ***  Виправлення Таврова Д.Ю. на основі файлу Яковлєва С.В. **

\documentclass{labs}

% Додаткові пакети вносьте у цей файл
%%%% У даний файл додавайте всі необхідні вам додаткові пакети, наприклад...


%%%% Діаграми
%\usepackage{tikz}                      % !!! невідомий конфлікт з якимось іншим пакетом

% Пакети для кольорових текстів (необхідні для команди \todo)
%\usepackage{xcolor}                     % !!! невідомий конфлікт з якимось іншим пакетом
%\usepackage{colortbl}

\usepackage{euscript}   % для красивих математичиних шрифтів \EuScript

%%%% ...і таке інше

%\usepackage[normalem]{ulem} % для підкреслень uline
%\ULdepth = 0.16em % визначаємо відстань від лінії до тексту вище/нижче

\usepackage{xstring} % For string manipulation


% Додаткові визначення та перевизначення команд вносьте у цей файл
%%% У даному файлі визначайте всі необхідні вам нові команди TeX
%%% або робіть перевизначення існуючих, наприклад...

% Перевизначення символу порожньої множини та знаків "більше-дорівнює", "менше-дорівнює" на прийняті у нас
\let\oldemptyset\emptyset
\let\emptyset\varnothing
\let\geq\geqslant
\let\leq\leqslant

% Визначення нових математичних команд
\newcommand*{\binsp}[1]{\ensuremath \left\{0, 1\right\}^{#1}}       % {0, 1}^m
\newcommand*{\xor}{\ensuremath \oplus}                              % \xor = (+)
\newcommand*{\GF}[1]{\ensuremath \mathbb F_{#1}}                    % F_n
\newcommand*{\GFgroup}[1]{\ensuremath \mathbb F^{*}_{#1}}           % F^*_n
\newcommand*{\Zring}[1]{\ensuremath \mathbb Z_{#1}}                 % Z_n
\newcommand*{\Zgroup}[1]{\ensuremath \mathbb Z^{*}_{#1}}            % Z^*_n
\newcommand*{\Jset}[1]{\ensuremath \mathbb J_{#1}}                  % J_n
\newcommand*{\Qset}[1]{\ensuremath \mathbb Q_{#1}}                  % Q_n
\newcommand*{\PQset}[1]{\ensuremath \widetilde{\mathbb Q}_{#1}}     % Q~_n
\newcommand*{\cyclic}[1]{\ensuremath \left\langle {#1} \right\rangle}                  % <g>
\newcommand*{\Legendre}[2]{\ensuremath \left( \frac{#1}{#2} \right)}  % символ Лежандра/Якоби
\newcommand*{\compinv}[1]{\ensuremath {#1}^{\left\langle -1 \right\rangle}}  % обратный по композиции

% Інший спосіб визначення математичного оператору
\DeclareMathOperator{\ord}{ord}
\DeclareMathOperator{\lcm}{lcm}
\DeclareMathOperator{\Li}{Li}
\DeclareMathOperator{\Coef}{Coef}
\DeclareMathOperator{\Log}{Log}
\DeclareMathOperator{\Exp}{Exp}
\DeclareMathOperator{\Res}{Res}
\DeclareMathOperator{\charact}{char}
\DeclareMathOperator{\Sym}{Sym}


% команда для коментарів червоним кольором
% !!! Конфлікт пакету color з якимось іншим пакетом, не використовувати
%\newcommand{\todo}[1]{\textcolor{red}{#1}}


%%% ...і таке інше

% Відомості про автора роботи
%%% Основні відомості %%%
%\newcommand{\UDC}                      % УДК
%{(впишіть правильний УДК!)}            % УДК виглядає приблизно як 004.056.5 або 513.2, або навіть 004.056.5:513.2+519.1
% Для того, щоб знайти правильний УДК, використовуйте каталог https://teacode.com/online/udc/

\newcommand{\reportAuthor}             % ПІБ автора повністю
{Іванов Петро Сидорович}
\newcommand{\reportAuthorGen}          % ПІБ автора в родовому відмінку
{Іванову Петру Сидоровичу}
\newcommand{\reportAuthorShort}        % ПІБ автора коротко
{Петро ІВАНОВ}
\newcommand{\reportAuthorGroup}        % група автора
{КМ-N3}
\newcommand{\reportTitle}              % Назва роботи
{Назва роботи}
%% використовуйте символ "\par" або "\\" для розбиття назви на декілька рядків

\newcommand{\supervisorFio}            % Науковий керівник, ПІБ повністю
{Прізвище Ім'я По-батькові}
\newcommand{\supervisorFioShort}       % Науковий керівник, ПІБ коротко
{Ім'я ПРІЗВИЩЕ}
\newcommand{\supervisorRegalia}        % Науковий керівник: посада, степінь, звання
{посада, ступінь, звання}              % наприклад: доцент кафедри ПМА, канд. техн. наук, доцент

\newcommand{\reportOrder}{\quotes{23} квітня 2024 р. \No ????-С} % реквізити наказу

\newcommand{\applicationDate}          % термін подання роботи
{\quotes{10} червня 2024 р.}

\newcommand{\assignmentDate}           % дата видачі завдання
{\quotes{05} лютого 2024 р.}

\newcommand{\consultFio}               % Консультант, ПІБ повністю
{}
\newcommand{\consultRegalia}           % Консультант: звання, ступінь, посада
{}
% Якщо у вас нема консультанта - залишайте ці поля порожніми


\newcommand{\reviewerFio}              % Рецензент, ПІБ повністю
{Прізвище Ім'я По-батькові}                        
\newcommand{\reviewerRegalia}          % Рецензент: звання, ступінь, посада
{посада, ступінь, звання}

\newcommand{\YearOfDefence}            % рік захисту
{2024}
\newcommand{\YearOfBeginning}          
{2023}

% Це файл з інформацією про джерела, на які посилаються в роботі
% Якщо формувати бібліографію вручну, то це варто закоментувати
\addbibresource{main.bib}

% Починаємо верстку документа
\begin{document}

% Створюємо титульну сторінку
\thispagestyle{empty}
\setlength{\parindent}{0cm}

\begin{center}
\textbf{\MakeUppercase{Національний технічний університет України}}

\textbf{\MakeUppercase{\quotes{Київський політехнічний інститут\\ імені Ігоря Сікорського}}}

\textbf{Факультет прикладної математики}

\textbf{Кафедра прикладної математики}

\vspace{5mm}

\begin{flushright}
	\renewcommand{\arraystretch}{0.8}
	\begin{tabular}{l}
		
		\quotes{До захисту допущено} \\
		Завідувач кафедри\\
		\_\_\_\_\_\_\_\_\_\_\_\_\_~Олег ЧЕРТОВ \\
		\quotes{\_\_\_\_}~\_\_\_\_\_\_\_\_\_\_\_\_\_~\the\year~р. \\
		
	\end{tabular}
\end{flushright}

\vfill 

\linespread{1.2}
\setfontsize{20}
\textbf{{Дипломна робота}}
\setfontsize{14}

\textbf{на здобуття ступеня бакалавра}

\textbf{за освітньо-професійною програмою \quotes{Наука про дані та математичне моделювання}}

\textbf{спеціальності 113 \quotes{Прикладна математика}}

\textbf{на тему: \quotes{\reportTitle}}
\linespread{1.5}

\end{center}

\vfill % невеликий відступ

{\renewcommand{\arraystretch}{0.7}
\begin{tabularx}{\textwidth}{>{\setlength\hsize{1.5\hsize}}X >{\setlength\hsize{0.5\hsize}}X}
Виконав:                                                 & \\
студент IV курсу, групи \reportAuthorGroup               & \\
\reportAuthor                                            & \rule{2.5cm}{0.25pt}   \\[6pt]
Керівник:                                                & \\
\supervisorRegalia                                       & \\
\supervisorFio                                           & \rule{2.5cm}{0.25pt}   \\[6pt]
%%%%% Якщо у вас є консультант у роботі - розкоментуйте наступні три рядки (а цей - не розкоментовуйте!)
%Консультант:                                             & \\
%\consultRegalia                                          & \\
%\consultFio                                              & \rule{2.5cm}{0.25pt}   \\[6pt]
Рецензент:                                               & \\
\reviewerRegalia                                         & \\
\reviewerFio                                             & \rule{2.5cm}{0.25pt} 
\end{tabularx}
}

\vfill % автоматично заповнюємо проміжок

\begin{flushright}
	\renewcommand{\arraystretch}{0.8}
	\begin{tabular}{l}
		Засвідчую, що в цій дипломній роботі\\
		немає запозичень із праць інших авторів\\
		без відповідних посилань.\\
		Студент~\_\_\_\_\_\_\_\_\_\_\_\_\_
	\end{tabular}
\end{flushright}

\vfill
\begin{center}
{Київ~---~\YearOfDefence\ року}
\end{center}

\clearpage % починаємо з нової сторінки

\setcounter{page}{2} % номер сторінки буде 2

\setlength{\parindent}{1.25cm} % задаємо відступ першого рядка абзацу (5 знаків, або 1,25 см)

% пропуск рядка до та після формули
\setlength{\belowdisplayskip}{14pt}
\setlength{\abovedisplayskip}{14pt}
\setlength{\belowdisplayshortskip}{14pt}
\setlength{\abovedisplayshortskip}{14pt}


% Створюємо зміст
%\pagenumbering{gobble}
\tableofcontents
\cleardoublepage
%\pagenumbering{arabic}
%\setcounter{page}{8}     %!!! -- продумати, як автоматизувати номер сторінки

% Створюємо вступ
\intro
%!TEX root = ../thesis.tex
% створюємо вступ
Актуальність даного дослідження полягає 
у тому, що без нього ви не одержите диплом про вищу освіту. Відповідно, ви повинні 
оформити результати вашого дослідження належним чином.



% Додаємо глави
%!TEX root = ../thesis.tex

\chapter{Постановка задачі}
\label{chap:problem}  %% відмічайте кожен розділ певною міткою -- на неї наприкінці необхідно посилатись

Першим розділом будь-якої роботи бакалавра має бути розділ про постановку задачі.

Тут треба вказати, серед іншого, мету і завдання.

\textit{Метою дослідження} є певна абстрактна недосяжна річ на кшталт 
загальнолюдського щастя на горизонті. Для досягнення мети потрібно 
розв'язати \textit{задачу дослідження}, яка полягає у чомусь суттєво більш 
конкретному. Для розв'язання задачі необхідно вирішити такі завдання:
\begin{enumerate}
	\item провести огляд опублікованих джерел за тематикою дослідження;
	\item (наступний пункт, пов'язаний із теоретичним дослідженням);
	\item (і ще один, наприклад, про експериментальну перевірку результатів);
	\item (і взагалі, краще з керівником проконсультуйтесь, як ваші 
	завдання правильно писати).
\end{enumerate}

Також у цьому розділі варто викласти вимоги до розроблюваного програмного забезпечення (як функціональні, так і нефункціональні, наприклад час виконання чи точність прогнозування тощо).
%!TEX root = ../thesis.tex

\chapter{(Назва другого розділу)}
\label{chap:review}

\section{(Назва першого підрозділу)}

Другий розділ, як правило, повинен бути присвячений огляду попередніх результатів за 
тематикою вашого дослідження. У даному розділі повинні міститись всі 
визначення та описи, потрібні для дальшого викладення матеріалу, та результати 
ваших попередників.

Абсолютно на всі не ваші результати повинні стояти 
в належний спосіб оформлені посилання.

Розмір другого (оглядового) розділу не повинен перевищувати третини вашої 
дипломної роботи (без урахування додатків).


\section{(Назва другого підрозділу)}

Наведемо основні правила оформлення текстів у системі \LaTeX.

Для абзацу робіть пусті рядки у файлі. Курсивний текст робиться командою 
\texttt{textit}: \textit{ось так}. 

\quotes{Лапки} робляться командою \texttt{quotes}. Довге 
тире у тексті --- трьома дефісами, коротке -- двома дефісами; у формулах 
мінуси робляться одним дефісом: $a-b$.

Пишіть звичайний текст звичайним текстом, а формули, позначення змінних та 
операцій (усі формули, усі позначення змінних та усі операції) беріть у 
знаки долара, ось так: $E = mc^2$, $a_1 = a^{(2)} \cdot a_{n, k}$, $e^x = 
\sum_{k = 0}^{\infty} {\frac{x^k}{k!}}$. Якщо вам 
не подобається, як \LaTeX подав формулу для експоненти (мені, наприклад, 
не подобається), то можна внести у код формули деякі корективи та написати ось так: $e^x 
= \sum\limits_{k = 0}^{\infty} {\dfrac{x^k}{k!}}$.

Для прикладу різні варіації коми у формулах: $(a, b)$ vs. $(a,b)$. Поки 
пакет \texttt{icomma} працює, різниця видна наочно.

Виключна формула (формула окремим рядком) робиться через спеціальне оточення з бекслешів та квадратових дужок або через оточення \texttt{equation}. Зауважте, що при цьому змінюється 
оформлення формул:
\[
e^x = \sum_{k = 0}^{\infty} {\frac{x^k}{k!}}\;.
\]
Оточення \texttt{equation} потрібне для створення нумерованих формул:
\begin{equation}
e^x = \sum_{k = 0}^{\infty} {\frac{x^k}{k!}}\;.
\label{eq:exponent}
\end{equation}

Потім на формули можна посилатися, наприклад, в \eqref{eq:exponent} наведено розвинення експоненти в ряд Маклорена.

Формули за замовчуванням не підтримують кирилічні літери. Зверніть увагу на 
порожній рядок перед попереднім реченням у tex-файлі: без нього не буде 
створено абзац.

Із більш специфічних позначень --- ось так, скажімо, можна подати 
перестановку:
$$\pi = \begin{pmatrix}
	1 & 2 & 3 & 4 & 5 & 6 & 7 & 8 & 9\\
	a & 5 & 9 & 6 & 4 & 8 & 2 & 1 & 7
\end{pmatrix},$$
де $a=3$. Зауважте, що у попередньому реченні нема порожнього рядочку 
перед <<де>> (та, відповідно, абзацу після формули), а кома внесена у 
виключну формулу, бо інакше вона переїде у наступний рядок тексту.

Декілька формул поспіль треба збирати в єдине ціле оточенням \texttt{align}; назви оточень із зірочками вказують \LaTeX'у не нумерувати формули. Наприклад, ось рекуренти для циклових чисел та чисел Стірлінга 
I~роду:
\begin{align*}
	c(n+1, k) &= c(n, k-1)+nc(n, k); \\
	s(n+1, k) &= s(n, k-1)-ns(n, k).
\end{align*}

Зверніть увагу на символ \verb|~| у попередньому абзаці tex-файлу; це нерозривний пробіл, який не дасть рознести пов'язані 
частини по різних рядках. Тільду треба ставити перед усіма посиланнями 
(команди ref та cite), перед тире та у місцях, які не можна розривати за 
правилами граматики.

Для специфічних позначень ви можете задавати власні команди (їх 
рекомендовано заносити у файл \texttt{02\_redefinitions}). Наприклад, 
подивіться, як оформлюється теорема Лагранжа-Бюрмана із використанням 
введених команд \texttt{Coef} та \texttt{compinv}:

\begin{theorem}[Лагранж, Бюрман] \label{thLagrangeBurmann}
	Для будь-якого ряду $A \in x \mathcal R[[x]]_1$ та $k \in \mathbb N$ справедливе співвідношення
	$$n \Coef[x^n] \left( \compinv{A}(x) \right)^k = k \Coef[x^{n-k}] \left(\! \frac{x}{A(x)} \!\right)^n.$$
\end{theorem}
\begin{proof}
	Доведення ви подивитесь деінде, а тут подивіться, як воно оформлюється 
	(зокрема, на квадратик наприкінці). 
\end{proof}

Іноді написаний файл треба компілювати двічі для одержання ефекту 
(скажімо, для коректної побудови усіх гіперпосилань та побудови змісту).  
Онлайн-сервіси на кшталт Overleaf справляються з такими ситуаціями за одну компіляцію. Однак той 
же Overleaf має звичку компілювати pdf-файли навіть за наявності помилок у 
тексті, просто ігноруючи відповідні місця. Якщо ви працюєте у Overleaf, 
то переконайтесь, що у вас нема червоних помилок після компіляції.

\section{(Назва третього підрозділу)}

Надамо деякі рекомендації щодо використання даного стильового файлу.

\begin{theorem}
	Використовуйте оточення \texttt{theorem} для теорем.
\end{theorem}
\begin{proof}
	Для доведень використовуйте оточення \texttt{proof}.
\end{proof}
\begin{theorem}
	Нумерація відбувається автоматично.
\end{theorem}
\begin{claim}
	Використовуйте оточення \texttt{claim} для тверджень.
\end{claim}
\begin{lemma}
	Використовуйте оточення \texttt{lemma} для лем.
\end{lemma}
\begin{corollary}
	Використовуйте оточення \texttt{corollary} для наслідків.
\end{corollary}
\begin{definition}
	Використовуйте оточення \texttt{definition} для визначень.
\end{definition}
\begin{example}
	Використовуйте оточення \texttt{example} для прикладів, на які є посилання.
\end{example}
\begin{remark}
	Використовуйте оточення \texttt{remark} для зауважень.
\end{remark}


\chapconclude

Наприкінці кожного розділу ви повинні навести коротенькі підсумки по його 
результатах. Зокрема, для оглядового розділу в якості висновків потрібно 
зазначити, які задачі у даній тематиці вже були розв'язані, а саме 
поставлена вами задача розв'язана не була (або розв'язана погано), тому у 
наступних розділах ви її й розв'язуєте.
%!TEX root = ../thesis.tex
% створюємо розділ
\chapter{(Назва третього розділу)}
\label{chap:theory}

\section{(Якийсь підрозділ)}

У третьому розділі варто наводити розв'язання поставленої перед вами 
задачі у теоретичному або аналітичному сенсі (хоча, звісно, все залежить 
від того, яка саме задача перед вами поставлена).

Для подання матеріалів можна використовувати таблиці (наприклад, 
табл. \ref{tab_weight}). 

    \begin{table}[ht]
    \caption{Розрахунок якоїсь величини у декілька кроків}
    \label{tab_weight}
    \centering
        \begin{tabular}{|c|c|c|c|c|c|c|c|c|}
        \hline \multirow{2}{*}{Параметр $x_i$} & \multicolumn{4}{c|}{Параметр $x_j$} & 
            \multicolumn{2}{c|}{Перший крок} & \multicolumn{2}{c|}{Другий крок} \\
        \cline{2-9} & $X_1$ & $X_2$ & $X_3$ & $X_4$ & $w_i$ & 
            ${K_\text{в}}_i$ & $w_i$ & ${K_\text{в}}_i$ \\
        \hline $X_1$ & 1 & 1 & 1.5 & 1.5 & 5 & 0.31 & 19 & 0.32 \\
        \hline $X_2$ & 1 & 1 & 1.5 & 1.5 & 5 & 0.31 & 19 & 0.32 \\
        \hline $X_3$ & 0.5 & 0.5 & 1 & 0.5 & 2.5 & 0.16 & 9.25 & 0.16 \\
        \hline $X_4$ & 0.5 & 0.5 & 1.5 & 1 & 3.5 & 0.22 & 12.25 & 0.20 \\
        \hline \multicolumn{5}{|c|}{Разом:} & 16 & 1 & 59.5 & 1 \\
        \hline
        \end{tabular}
    \end{table}

Бажано, щоб кожен пункт завдань, окреслених у вступі, відповідав певному 
розділу або підрозділу у дипломній роботі.

\begin{theorem}
Нумерація у наступних розділах також проставляється автоматично та коректно.
\end{theorem}

\section{(Якийсь наступний підрозділ)}

Для подання матеріалів також дуже зручними є рисунки (наприклад, рис.
\ref{fig_sudak}).


\begin{figure}[ht]
\centering
    \begin{subfigure}[b]{0.5\textwidth}    
        \includegraphics[scale=0.3]{Images/Sudak.png}
        \caption{}
    \end{subfigure}%
    \begin{subfigure}[b]{0.5\textwidth}
        \includegraphics[scale=0.3]{Images/Tudak.png}
        \caption{}
    \end{subfigure}
 
    \caption{Різні види риб: (a) судак, (б) тудак}
    \label{fig_sudak}
\end{figure}

\chapconclude

Наприкінці розділу знову наводяться коротенькі підсумки.

% Створюємо висновки
\conclusions
%!TEX root = ../thesis.tex
% створюємо Висновки до всієї роботи
Загальні висновки до роботи повинні підсумовувати усі ваші досягнення у 
даному напрямку досліджень.

За кожним пунктом завдань, поставлених у вступі, у висновках повинен 
міститись звіт про виконання: виконано, не виконано, виконано частково (і 
чому саме так). Наприклад, якщо першим поставленим завданням у вас іде 
\quotes{огляд літератури за тематикою досліджень}, то на початку висновків ви 
повинні зазначити, що \quotes{у ході роботи був проведений аналіз 
опублікованих джерел за тематикою (...), який показав, що (...)}. Окрім 
простої констатації про виконання ви повинні навести, які саме результати 
ви одержали та проінтерпретувати їх з точки зору поставленої задачі, мети 
та загальної проблематики.

В ідеалі загальні висновки повинні збиратись з висновків до кожного 
розділу. Однак висновки не повинні містити 
формул, таблиць та рисунків. У висновках обов'язково повинні фігурувати конкретні числа
(на кшталт \quotes{нейронна мережа прогнозує з точністю $95.71\%$}).

Наприкінці висновків потрібно зазначити напрямки подальших досліджень: 
куди саме, як вам вважається, потрібно прямувати наступним дослідникам у 
даній тематиці.


% Додаємо бібліографію
%%%%%%%% Якщо ви використовуєте bibtex, оце все треба розкоментувати звідси...
\defbibenvironment{bibliography}
  {\list
     {\printfield[labelnumberwidth]{labelnumber}}
     {\setfontsize{14}%
     }}
  {\endlist}
  {\item}

\printbibliography[heading=bibintoc,title={Перелік посилань}]
%%%%%%%% ...по осюди

% Замість бібтеху можна вписати та оформити бібліографію вручну
% %!TEX root = ../thesis.tex
% створюємо список використаної літератури
\begin{thebibliography}    
    \bibitem{Chaffey}
    Chaffey D. Digital marketing: strategy, implemenation and practice / Dave Chaffey, Ellis-Chadwick Fiona. --- 6th ed. --- [S. l.] : Pearson, 2016. --- 702 p.
    
    \bibitem{Daniels}
    Daniels D. Email Marketing: An Hour a Day / David Daniels, Jeanniey Mullen. Indianapolis : Wiley Publishing, 2009. --- 291 p.
    
    \bibitem{Comrack}
    Comrack G. V. Email Spam Filtering A Systematic Review / Gordon V. Comrack. --- Hanover : Now Publishers, 2008. --- 136 p.
    
    \bibitem{manage}
    Manage domains --- Postmaster tools [Electronic resource] // Google Postmaster tools. --- Mode of access: \url{https://postmaster.google.com} (date of access: 29.04.2024). --- Title from screen.
    
    \bibitem{Wolford}
    Wolford B. How does the GDPR affect email? [Electronic resource] / Ben Wolford // GDPR.EU. --- Mode of access: \url{https://gdpr.eu/email-encryption/} (date of access: 28.04.2024). --- Title from screen.
    
    \bibitem{privacy}
    Privacy and Identity Management [Electronic resource] / ed. by M. Friedewald, S. Schiffner, S. Krenn. --- Cham : Springer International Publishing, 2021. --- Mode of access: \url{https://doi.org/10.1007/978-3-030-72465-8} (date of access: 28.04.2024). --- Title from screen.
    
    \bibitem{turn}
    Turn Emails into Revenue [Electronic resource] // Intuit Mailchimp. --- Mode of access: \url{https://mailchimp.com} (date of access: 11.05.2024). --- Title from screen.

 
\end{thebibliography}


% Створюємо додатки (дивись у файли додатків для пояснень)
% Якщо ви маєте меншу або більшу кількість додатків, модифікуйте ці рядки у відповідний спосіб
% Якщо ви не маєте додатків, просто закоментуйте ці рядки:
%!TEX root = ../thesis.tex
\append{Лістинги програм}
\label{appendix:A}

\section{Програма 1}

Лістинг файлу main.py

\lstinputlisting[language=Python]{Code/main.py}



Зауважте, як змінилась нумерація підрозділу.


% Нарешті
\end{document}
